\subsection{The geometric recognition maps}
\label{subsec:geometric-recognition-maps}

This subsection composes the line-tangency lemma with the
covering and monodromy maps for the punctured elliptic curve.  We arrange that the
detecting compactification has genus at least ten, so the published
stable-range inputs apply directly.
Let \(X_0=E_0-\{O\}\), where \(E_0:y^2=x^3-x\) over \(\mathbb Q\),
and choose a nonzero rational tangent vector at \(O\).
Write
\[
 F=\pi_1(X_{0,\overline{\mathbb Q}})
       =\widehat{\langle x,y\rangle},\qquad
 \Pi=F\rtimes s_O(G_{\mathbb Q}).
\]
The cusp inertia \(I_O\) is generated by \([x,y]\), up to inversion.
All subgroups and commutator closures below are closed.

\paragraph{A fixed geometric cover in the stable range.}
Consider the explicit surjection
\[
 F\longrightarrow \operatorname{UT}_3(\mathbb F_3),\qquad
 x\longmapsto1+E_{12},\qquad y\longmapsto1+E_{23}.
\]
The two matrices and their commutator generate the group of order
\(27\), and the cusp generator maps to \(1+E_{13}\), of order three.
Let \(F_{10}\) be its kernel and \(Y_{10}\to X_{0,\overline{\mathbb Q}}\)
the corresponding connected finite \'etale cover.
Its compactification \(C_{10}\) has nine points over \(O\),
each of ramification index three.  Riemann--Hurwitz gives
\[
 2g(C_{10})-2=27(2g(E_0)-2)+9(3-1)=18.
\]
Thus \(g(C_{10})=10\).
The cover and its morphism descend to a finite number field,
since they are given by finitely many algebraic equations.
Only this actual geometric covering subgroup will be used.

\begin{theorem}\label{thm:geometric-procyclic-recognition}
Every nontrivial closed procyclic subgroup of \(F\) whose
arithmetic normalizer has open image in \(G_{\mathbb Q(i)}\)
is contained in a conjugate of \(I_O\).
The same conclusion holds after every finite arithmetic base
extension.  Conversely, every closed subgroup of a cusp inertia
group has an arithmetic normalizer with open image.
\end{theorem}

\begin{proof}
Let \(A=\overline{\langle z\rangle}\neq1\) and let \(N\) be its
arithmetic normalizer.  Its image is open in \(G_{\mathbb Q}\),
because \(G_{\mathbb Q(i)}\) is open.
Suppose that \(z\) is outside the union \(\mathcal Z\) of the
conjugates of cusp inertia.
Nakamura's full-class proper-homology detector
\cite[Lemma 2.1.4, pp.~96--97]{Nakamura1994}, applied before any
pro-prime quotient, gives a prime \(\ell\) and an open
\(H\subset F\) containing \(z\) such that
\begin{equation}\label{eq:geo-proper-detector}
 z\notin[H,H]H^\ell\langle\mathcal Z\cap H\rangle.
\end{equation}
Let \(C_H\) compactify the finite geometric cover defined by \(H\).
Killing the boundary inertia is the actual filling homomorphism
\(H\twoheadrightarrow\pi_1(C_H)\), so the quotient in
\eqref{eq:geo-proper-detector} is
\(H_1(C_H,\mathbb F_\ell)\).
Consequently the specified element \(z\) has a nonzero class
in the free \(\mathbb Z_\ell\)-module \(H_1(C_H,\mathbb Z_\ell)\).

The homomorphism to this module kills the primary factors
of \(A\) at primes different from \(\ell\).
If \(u\) is the \(\ell\)-primary component of \(z\), then
\(A_\ell=\overline{\langle u\rangle}\cong\mathbb Z_\ell\)
and \([u]\neq0\) in \(H_1(C_H,\mathbb Q_\ell)\).
This primary factor is characteristic in \(A\), hence is
normalized by \(N\).

Form the full arithmetic core
\begin{equation}\label{eq:geo-arithmetic-core}
 S=\bigcap_{\gamma\in\Pi}\gamma(H\cap F_{10})\gamma^{-1}.
\end{equation}
It is open in \(F\): a finitely generated profinite group
has only finitely many subgroups of any specified finite index,
as follows from their finite permutation representations.
Thus the intersection in \eqref{eq:geo-arithmetic-core} is finite.
It is normal in \(\Pi\) and contained in \(H\cap F_{10}\).
For some \(r\geq0\), the element \(u'=u^{\ell^r}\) generates
\(A_\ell\cap S\).
The inclusions of covering groups and the associated maps of
proper compactifications give the exact pushforward
\begin{equation}\label{eq:geo-pushforward-class}
 H_1(C_S,\mathbb Q_\ell)\longrightarrow H_1(C_H,\mathbb Q_\ell),
 \qquad [u']\longmapsto\ell^r[u]\neq0.
\end{equation}
The map \(C_S\to C_{10}\) is finite and nonconstant.
Riemann--Hurwitz therefore gives \(g(C_S)\geq10\).

The arithmetic subgroup
\[
                  \Pi_Y=S\rtimes s_O(G_{\mathbb Q})
\]
defines an actual geometrically connected cover \(Y/\mathbb Q\)
of \(X_0\).  At the chosen branch over \(O\), its decomposition
subgroup is exactly
\[
 \Pi_Y\cap\bigl(I_O\rtimes s_O(G_{\mathbb Q})\bigr)
               =(S\cap I_O)\rtimes s_O(G_{\mathbb Q}).
\]
Its residue Galois image is all of \(G_{\mathbb Q}\).
Hence that branch is a rational point \(y\) on the smooth
proper compactification \(C=\overline Y\).
This is an explicit descent through an arithmetic covering
subgroup; it does not use a field-of-moduli assertion.

The open immersion
\[
                  Y\longrightarrow X=C-\{y\}
\]
fills the other punctures.  It induces a surjection of geometric
fundamental groups and the identity on the arithmetic quotient.
On proper homology it is the identity.
Consequently the image of \(u'\) in
\[
                         P=\pi_1(X_{\overline{\mathbb Q}})^{(\ell)}
\]
has a nonzero class \(a\) in
\(P^{\mathrm{ab}}\otimes\mathbb Q_\ell
 =H_1(C,\mathbb Q_\ell)=:V\).
In particular the image of \(A_\ell\cap S\) is still
isomorphic to \(\mathbb Z_\ell\): a nonzero closed kernel
would give a finite image, whose rational abelian class is zero.
The image of \(N\cap\Pi_Y\) is a closed normalizer with the
same open arithmetic image.
Choose a finite number field \(L\) such that this image
contains \(G_L\).

\paragraph{The actual degree-six arithmetic quotient.}
For the same curve \(X\), Nakamura
\cite[Corollary 4.2 and the table in 4.3, p.~115;
Lemma 4.5, p.~116; (4.7.1)--(4.7.2), p.~117]{Nakamura1996}
provides the following exact maps.
The universal outer arithmetic degree-six image \(U_{g,6}\)
has a quotient
\[
                  q_{g,6}:U_{g,6}\twoheadrightarrow Q_{g,6},
 \qquad \dim_{\mathbb Q_\ell}Q_{g,6}=1.
\]
The action on the quotient is the cube of the symplectic
similitude character.
The actual degree-six graded Galois image of \(X/L\)
surjects onto this quotient.
The degree-six dimension follows from the rank sandwich
whose two endpoints are both one; the source also states
finite-base invariance immediately before Lemma 4.5.
After clearing a rational denominator, there is therefore
an actual
\[
 \sigma\in G_L(6),\qquad D_\sigma\in U_{g,6},\qquad
                              q_{g,6}(D_\sigma)\neq0.
\]
Here \(D_\sigma\) is the canonical special representative
of its outer graded class.

Lift this same \(\sigma\) to the normalizer of the image
of \(A_\ell\cap S\).  Since \(\sigma\) acts trivially on \(V\),
its normalizer exponent \(\lambda\in\mathbb Z_\ell^\times\)
satisfies \(\lambda a=a\), hence \(\lambda=1\).
Write \(U=\log u'=a+\text{terms of degree at least two}\)
in the rational completed free Lie algebra.
The chosen normalizing automorphism and the canonical
representative differ by an inner automorphism, so
\[
                 \phi_\sigma(U)=\exp(\operatorname{ad}Z)U.
\]
Its degree-one conjugator term is a multiple of \(a\).
Multiplying the conjugator by the inverse of the appropriate
rational power of \(\exp U\) removes that term and preserves
its action on \(U\).
The degree-two through degree-five terms then centralize \(a\)
and vanish, because the free-Lie centralizer of a nonzero
free generator is its one-dimensional span.
Taking degree seven gives the exact equality
\begin{equation}\label{eq:geo-normalizer-tangency}
                         D_\sigma(a)=[Z_6,a].
\end{equation}
Using the full \(U\), rather than only its linear term,
is necessary for this conjugator correction.

\paragraph{The geometric image in the certificate.}
We verify the monodromy comparison needed to apply the
polynomial identity to this \(D_\sigma\).
Use the moduli stack with one marked point and a nonzero
tangent vector.  Its geometric mapping class group fixes
the boundary word and has Zariski closure \(R_g\) in
the automorphisms of the rational punctured Malcev group.
At each finite nilpotent level, the actual profinite
geometric image has the same Zariski closure after
extension to \(\mathbb Q_\ell\): the discrete image is
dense and the defining polynomials are continuous.

Hain's Torelli-to-relative-completion surjection
\cite[Theorem 3.4, p.~605]{Hain1997} identifies the
prounipotent radical mapping onto \((R_g)_u\).
His weight and boundary-action comparison
\cite[Theorem 4.10, p.~611; Remark 4.11, p.~613;
\S14.4, p.~643]{Hain1997} is a morphism of mixed
Hodge structures and is strict.
Consequently its graded image is generated by its
degree-one image.  Enomoto--Satoh
\cite[Theorem 3.5 and Proposition 3.6, p.~640]{EnomotoSatoh2014}
identifies that generated image with the rational Johnson
image \(J_m\).  In particular,
\[
              \operatorname{gr}_{-m}\operatorname{Lie}(R_g)_u=J_m,
              \qquad J_1=\bigwedge^3V.
\]
This proves the equality for the image filtration; density
alone is not used to identify individual filtration kernels.
In degree six, the special-to-outer comparison is injective:
an inner derivation preserving the symplectic boundary class
has generator in its free-Lie centralizer, which is the
one-dimensional span of that degree-two class.
Thus the geometric kernel in Nakamura's exact sequence is
the same \(J_6\), and \(q_{g,6}\) kills the actual Johnson
brackets used in the certificate.

Arithmetic monodromy normalizes the geometric algebraic image.
For an outer degree-six class its tangential representative
can have a lower inner term only from the degree-two
boundary class.  Subtracting this term by
\(\exp(-c\,\operatorname{ad}\log(\partial))\) stays inside
the geometric algebraic image: the boundary Dehn twist
generates its one-parameter Zariski closure.
The centralizer of the boundary class removes the
remaining lower inner terms.
The resulting canonical special representative therefore
still normalizes \(R_g\) to the required order.
Taking the leading term and using strictness gives
\[
                         [D_\sigma,J_1]\subseteq J_7.
\]

Thus the Enomoto--Satoh cyclic contraction and the additional
five-term contraction verified by the six Johnson vectors
both vanish on these brackets.  Their seven scalar rows
are precisely the necessary rows of the polynomial certificate.
All stable-range requirements hold because \(g\geq10\).
Choose \(\langle a,b\rangle=1\), set
\(W=\langle a,b\rangle^\perp\), and let
\(\operatorname{Sp}(W)\times\mathbb G_m\) act by its standard
action on \(W\) and by \(t,t^{-1}\) on \(a,b\).
Its reductive invariant projector preserves both the actual universal
image and the tangency condition
\eqref{eq:geo-normalizer-tangency}; it preserves \(q_{g,6}\),
because that quotient is symplectically trivial.
Lemma~\ref{lem:geometric-line-tangency} therefore proves
\[
 D\in U_{g,6},\quad D(a)\in[\mathbb L_6(V),a]
                       \quad\Longrightarrow\quad q_{g,6}(D)=0.
\]
Applied to \(D_\sigma\), this contradicts its nonzero quotient.
Hence \(z\in\mathcal Z\), and its procyclic closure lies
in that same cusp inertia group.

The argument depends only on an open arithmetic image.
Starting after a finite arithmetic base extension gives
the same geometric group and the same inertia conjugates;
the normalizer image remains open in \(G_{\mathbb Q}\).
Finally, a cusp decomposition group over a field defining
the cusp acts on its procyclic inertia by the cyclotomic
character.  Every closed subgroup of a procyclic group
is characteristic.  That decomposition group therefore
normalizes it and has open arithmetic image, proving
the converse.
\end{proof}

\begin{corollary}\label{cor:geometric-common-kernel-recognition}
Put
\[
 R_{\mathrm{geom}}=
 \bigcap_\ell\ker\!\left(F\longrightarrow F^{(\ell)}\right).
\]
No nontrivial closed procyclic subgroup of \(R_{\mathrm{geom}}\)
has an arithmetic normalizer with open image.
\end{corollary}

\begin{proof}
The intersection of each cusp inertia group with
\(R_{\mathrm{geom}}\) is trivial by an explicit finite quotient.
For every \(n\geq1\), the substitutions
\[
 x\longmapsto1+E_{12},\qquad y\longmapsto1+E_{23}
       \quad\text{in }\operatorname{UT}_3(\mathbb Z/\ell^n\mathbb Z)
\]
send the cusp commutator to \(1+E_{13}\), of exact order
\(\ell^n\).
Thus if a cusp element has exponent
\(\alpha=(\alpha_\ell)_\ell\in\widehat{\mathbb Z}\),
its pro-\(\ell\) image is nontrivial whenever
\(\alpha_\ell\neq0\).
Every nonzero exponent has such a component.
The equality persists under conjugation, since
\(R_{\mathrm{geom}}\) is normal.
The theorem now gives the corollary.
\end{proof}

The arithmetic outer kernel is related to the centralizer in the
fundamental exact sequence by the isomorphism proved in
Subsection~\ref{subsec:cm-geometric-common-kernel}.
The exact pullback calculation there identifies the extra
morphism needed to apply this corollary to that kernel.
